\ifdefined\gkver\else\def\gkver{student}\fi
\documentclass[\gkver]{gaokaozhenti}
\begin{document}

\chapter{2025年甘肃}

\begin{enumerate}
\item
%% number: 1
%% paperName: 2025年甘肃高考第 1 题 4 分
%% typeId: 1
%% type: 单选题
%% chapter: 原子和原子核
%% point: 玻尔模型
%% method: 无
%% score: 4
%% degree: 850
%% duplicateId: 0
%% body:
利用电子与离子的碰撞可以研究离子的能级结构和辐射特性。\ce{He+} 离子相对基态的能级图（设基态能量为 0）如图所示。用电子碰撞 \ce{He+} 离子使其从基态激发到可能的激发态，若所用电子的能量为 $50 \UeV$，则 \ce{He+} 离子辐射的光谱中，波长最长的谱线对应的跃迁为\xzanswer{C}

\onepicture[w=3.80cm]{figs/01.png}

\fourchoices[answer=C]
{$n = 4 \to n = 3$ 能级}
{$n = 4 \to n = 2$ 能级}
{$n = 3 \to n = 2$ 能级}
{$n = 3 \to n = 1$ 能级}

%% answer: C
%% memo:
\memoanswer{
根据题意可知，用能量为 $50 \UeV$ 的电子碰撞 \ce{He+} 离子，可使 \ce{He+} 离子跃迁到 $n = 3$ 能级和 $n = 2$ 能级，由

$\Delta E = E_{m} - E_{n} = h\nu = h\frac{c}{\lambda}$

可知，波长最长的谱线对应的跃迁为 $n = 3 \to n = 2$ 能级。

故选 C。
}

\item
%% number: 2
%% paperName: 2025年甘肃高考第 2 题 4 分
%% typeId: 1
%% type: 单选题
%% chapter: 曲线运动
%% point: 天体运动
%% method: 无
%% score: 4
%% degree: 750
%% duplicateId: 0
%% body:
如图，一小星球与某恒星中心距离为 $R$ 时，小星球的速度大小为 $v$、方向与两者中心连线垂直。恒星的质量为 $M$，引力常量为 $G$。下列说法正确的是\xzanswer{A}

\onepicture[w=3.40cm]{\includesvg{figs/02}}

\fourchoices[answer=A]
{若 $v = \sqrt{\frac{GM}{R}}$，小星球做匀速圆周运动}
{若 $\sqrt{\frac{GM}{R}} < v < \sqrt{\frac{2GM}{R}}$，小星球做抛物线运动}
{若 $v = \sqrt{\frac{2GM}{R}}$，小星球做椭圆运动}
{若 $v > \sqrt{\frac{2GM}{R}}$，小星球可能与恒星相撞}

%% answer: A
%% memo:
\memoanswer{
A．根据题意，由万有引力提供向心力有 $\frac{GMm}{R^{2}} = m\frac{v^{2}}{R}$

解得 $v = \sqrt{\frac{GM}{R}}$

可知，若 $v = \sqrt{\frac{GM}{R}}$，小星球做匀速圆周运动，故 A 正确；

B．结合 A 分析可知，若 $\sqrt{\frac{GM}{R}} < v < \sqrt{\frac{2GM}{R}}$，万有引力不足以提供小星球做匀速圆周运动所需要的向心力，小星球做离心运动，但又不能脱离恒星的引力范围，所以小星球做椭圆运动，而不是抛物线运动，故 B 错误；

C．若 $v = \sqrt{\frac{2GM}{R}}$，这是小星球脱离恒星引力束缚的临界速度，小星球将做抛物线运动，而不是椭圆运动，故 C 错误；

D．若 $v > \sqrt{\frac{2GM}{R}}$，小星球将脱离恒星引力束缚，做双曲线运动，不可能与恒星相撞，故 D 错误。

故选 A。
}

\item
%% number: 3
%% paperName: 2025年甘肃高考第 3 题 4 分
%% typeId: 1
%% type: 单选题
%% chapter: 运动和力
%% point: 牛顿第二定律
%% method: 无
%% score: 4
%% degree: 850
%% duplicateId: 0
%% body:
2025 年 4 月 24 日，在甘肃酒泉卫星发射中心成功发射了搭载神舟二十号载人飞船的长征二号 F 遥二十运载火箭。若在初始的 $1 \Us$ 内燃料对火箭的平均推力约为 $6 \times 10^{6} \UN$。火箭质量约为 $500 \Ut$ 且认为在 $1 \Us$ 内基本不变，则火箭在初始 $1 \Us$ 内的加速度大小约为\xzanswer{A}（重力加速度 $g$ 取 $10 \Umsq$）

\fourchoices[answer=A]
{$2 \Umsq$}
{$4 \Umsq$}
{$6 \Umsq$}
{$12 \Umsq$}

%% answer: A
%% memo:
\memoanswer{
根据题意，由牛顿第二定律有 $F - mg = ma$

代入数据解得 $a = \frac{6 \times 10^{6} - 5 \times 10^{6}}{5 \times 10^{5}} \Umsq = 2 \Umsq$

故选 A。
}

\item
%% number: 4
%% paperName: 2025年甘肃高考第 4 题 4 分
%% typeId: 1
%% type: 单选题
%% chapter: 运动和力
%% point: 动量和能量
%% method: 无
%% score: 4
%% degree: 650
%% duplicateId: 0
%% body:
如图，小球 A 从距离地面 $20 \Um$ 处自由下落，$1 \Us$ 末恰好被小球 B 从左侧水平击中，小球 A 落地时的水平位移为 $3 \Um$。两球质量相同，碰撞为完全弹性碰撞，重力加速度 $g$ 取 $10 \Umsq$，则碰撞前小球 B 的速度大小 $v$ 为\xzanswer{B}

\onepicture[w=5.50cm]{\includesvg{figs/04}}

\fourchoices[answer=B]
{$1.5 \Ums$}
{$3.0 \Ums$}
{$4.5 \Ums$}
{$6.0 \Ums$}

%% answer: B
%% memo:
\memoanswer{
根据题意可知，小球 A 和 B 碰撞过程中，水平方向上动量守恒，竖直方向上 A 球的竖直速度不变，设碰撞后 A 球水平速度为 $v_{1}$，B 球水平速度为 $v_{2}$，则有 $mv = mv_{1} + mv_{2}$

碰撞为完全弹性碰撞，则由能量守恒定律有 $\frac{1}{2}mv^{2} + \frac{1}{2}mv_{A}^{2} = \frac{1}{2}mv_{A}^{2} + \frac{1}{2}mv_{1}^{2} + \frac{1}{2}mv_{2}^{2}$

联立解得 $v_{1} = v$，$v_{2} = 0$

小球 A 在竖直方向上做匀加速直线运动，则有 $h = \frac{1}{2}gt^{2}$

解得 $t = 2 \Us$

可知，碰撞后，小球 A 运动 $t' = 1 \Us$ 落地，则水平方向上有 $x = vt'$

解得 $v = 3.0 \Ums$

故选 B。
}

\item
%% number: 5
%% paperName: 2025年甘肃高考第 5 题 4 分
%% typeId: 1
%% type: 单选题
%% chapter: 电场
%% point: 电势能电势差电场力做功
%% method: 无
%% score: 4
%% degree: 650
%% duplicateId: 0
%% body:
如图，两极板不平行的电容器与直流电源相连，极板间形成非匀强电场，实线为电场线，虚线表示等势面。M、N 点在同一等势面上，N、P 点在同一电场线上。下列说法正确的是\xzanswer{D}

\onepicture[w=4.30cm]{figs/05.png}

\fourchoices[answer=D]
{M 点的电势比 P 点的低}
{M 点的电场强度比 N 点的小}
{负电荷从 M 点运动到 P 点，速度增大}
{负电荷从 M 点运动到 P 点，电场力做负功}

%% answer: D
%% memo:
\memoanswer{
A．MN 两点电势相等，电场线由上到下，NP 在同一电场线上，沿电场线电势逐渐降低，可知 N 点电势高于 P 点，可知 M 点电势高于 P 点，选项 A 错误；

B．M 点电场线分布比 N 点密集，可知 M 点电场强度比 N 点大，选项 B 错误；

CD．负电荷从 M 点运动到 P 点，电势能增加，则电场力做负功，动能减小，速度减小，选项 C 错误，D 正确；

故选 D。
}

\item
%% number: 6
%% paperName: 2025年甘肃高考第 6 题 4 分
%% typeId: 1
%% type: 单选题
%% chapter: 电磁感应
%% point: 电磁感应中图象
%% method: 无
%% score: 4
%% degree: 550
%% duplicateId: 0
%% body:
闭合金属框放置在磁场中，金属框平面始终与磁感线垂直。如图，磁感应强度 $B$ 随时间 $t$ 按正弦规律变化。$\Phi$ 为穿过金属框的磁通量，$E$ 为金属框中的感应电动势，下列说法正确的是\xzanswer{C}

\onepicture[w=5.20cm]{\includesvg{figs/06}}

\fourchoices[answer=C]
{$t$ 在 $0 \sim \frac{T}{4}$ 内，$\Phi$ 和 $E$ 均随时间增大}
{当 $t = \frac{T}{8}$ 与 $\frac{3T}{8}$ 时，$E$ 大小相等，方向相同}
{当 $t = \frac{T}{4}$ 时，$\Phi$ 最大，$E$ 为零}
{当 $t = \frac{T}{2}$ 时，$\Phi$ 和 $E$ 均为零}

%% answer: C
%% memo:
\memoanswer{
A．在 $0 \sim \frac{T}{4}$ 间内，磁感应强度 $B$ 增加，根据 $\Phi = BS$ 则磁通量 $\Phi$ 增加，但是图像的斜率减小，即磁感应强度 $B$ 的变化率逐渐减小，根据法拉第电磁感应定律 $E = \frac{\Delta B}{\Delta t}S$ 可知，感应电动势 $E$ 逐渐减小，选项 A 错误；

B．当 $t = \frac{T}{8}$ 与 $\frac{3T}{8}$ 时，因 $B-t$ 图像的斜率大小相等，符号相反，可知感应电动势 $E$ 大小相等，方向相反，选项 B 错误；

C．$t = \frac{T}{4}$ 时，$B$ 最大，则磁通量 $\Phi$ 最大，但是 $B$ 的变化率为零，则感应电动势 $E$ 为零，选项 C 正确；

D．$t = \frac{T}{2}$ 时，$B$ 为零，则磁通量 $\Phi$ 为零，但是 $B$ 的变化率最大，则感应电动势 $E$ 最大，选项 D 错误。

故选 C。
}

\item
%% number: 7
%% paperName: 2025年甘肃高考第 7 题 4 分
%% typeId: 1
%% type: 单选题
%% chapter: 电场
%% point: 带电粒子的偏转
%% method: 无
%% score: 4
%% degree: 350
%% duplicateId: 0
%% body:
离子注入机是研究材料辐照效应的重要设备，其工作原理如图~\ref{2025甘肃07a}所示。从离子源 S 释放的正离子（初速度视为零）经电压为 $U_{1}$ 的电场加速后，沿 OO$'$ 方向射入电压为 $U_{2}$ 的电场（OO$'$ 为平行于两极板的中轴线）。极板长度为 $l$、间距为 $d$，$U_{2}-t$ 关系如图~\ref{2025甘肃07b}所示。长度为 $a$ 的样品垂直放置在距 $U_{2}$ 极板 $L$ 处，样品中心位于 O$'$ 点。假设单个离子在通过 $U_{2}$ 区域的极短时间内，电压 $U_{2}$ 可视为不变，当 $U_{2} = \pm U_{m}$ 时。离子恰好从两极板的边缘射出。不计重力及离子之间的相互作用。下列说法正确的是\xzanswer{B}

\twopicture[labelA=2025甘肃07a, labelB=2025甘肃07b, numA=1, numB=2, wA=8.40cm, wB=6.00cm]{figs/07a.png}{figs/07b.png}

\fourchoices[answer=B]
{$U_{2}$ 的最大值 $U_{m} = \frac{d^{2}}{l^{2}}U_{1}$}
{当 $U_{2} = \pm U_{m}$ 且 $L = \frac{(a - d)l}{2d}$ 时，离子恰好能打到样品边缘}
{若其他条件不变，要增大样品的辐照范围，需增大 $U_{1}$}
{在 $t_{1}$ 和 $t_{2}$ 时刻射入 $U_{2}$ 的离子，有可能分别打在 A 和 B 点}

%% answer: B
%% memo:
\memoanswer{
A．粒子在加速电场中被加速时

$U_{1}q = \frac{1}{2}mv_{0}^{2}$

在偏转电场中做类平抛运动，则 $l = v_{0}t$，$\frac{d}{2} = \frac{1}{2} \cdot \frac{U_{m}q}{dm}t^{2}$

解得 $U_{m} = \frac{2d^{2}}{l^{2}}U_{1}$

选项 A 错误；

B．当 $U_{2} = \pm U_{m}$ 时粒子从板的边缘射出，恰能打到样品边缘时，则

$\frac{\frac{d}{2}}{\frac{a}{2}} = \frac{\frac{l}{2}}{\frac{l}{2} + L}$

解得 $L = \frac{(a - d)l}{2d}$

选项 B 正确；

C．根据 $y = \frac{1}{2} \cdot \frac{U_{m}q}{dm}t^{2} = \frac{U_{m}l^{2}}{4dU_{1}}$

若其它条件不变，要增加样品的辐照范围，则需减小 $U_{1}$，选项 C 错误；

D．由图可知 $t_{1}$ 时刻所加的向上电场电压小于 $t_{2}$ 时刻所加的向下的电场的电压，则 $t_{1}$ 时刻射入的粒子打到 A 点时的竖直位移小于打到 B 点时的竖直位移，则选项 D 错误。故选 B。
}

\item
%% number: 8
%% paperName: 2025年甘肃高考第 8 题 5 分
%% typeId: 2
%% type: 多选题
%% chapter: 运动和力
%% point: 瞬时加速度
%% method: 无
%% score: 5
%% degree: 550
%% duplicateId: 0
%% body:
如图，轻质弹簧上端固定，下端悬挂质量为 $2m$ 的小球 A，质量为 $m$ 的小球 B 与 A 用细线相连，整个系统处于静止状态。弹簧劲度系数为 $k$，重力加速度为 $g$。现剪断细线，下列说法正确的是\xzanswer{BC}

\onepicture[h=5.80cm]{\includesvg{figs/08}}

\fourchoices[answer=BC]
{小球 A 运动到弹簧原长处的速度最大}
{剪断细线的瞬间，小球 A 的加速度大小为 $\frac{g}{2}$}
{小球 A 运动到最高点时，弹簧的伸长量为 $\frac{mg}{k}$}
{小球 A 运动到最低点时，弹簧的伸长量为 $\frac{2mg}{k}$}

%% answer: BC
%% memo:
\memoanswer{
A．剪断细线后，弹力大于 A 的重力，则 A 先向上做加速运动，随弹力的减小，则向上的加速度减小，当加速度为零时速度最大，此时弹力等于重力，弹簧处于拉伸状态，选项 A 错误；

B．剪断细线之前则 $F_{\text{弹}} = 3mg$

剪断细线瞬间弹簧弹力不变，则对 A 由牛顿第二定律 $F_{\text{弹}} - 2mg = 2ma$

解得 A 的加速度 $a = \frac{g}{2}$

选项 B 正确；

C．剪断细线之前弹簧伸长量 $x_{1} = \frac{3mg}{k}$

剪断细线后 A 做简谐振动，在平衡位置时弹簧伸长量 $x_{2} = \frac{2mg}{k}$

即振幅为 $A = x_{1} - x_{2} = \frac{mg}{k}$

由对称性可知小球 A 运动到最高点时，弹簧伸长量为 $\frac{mg}{k}$，选项 C 正确；

D．由上述分析可知，小球 A 运动到最低点时，弹簧伸长量为 $\frac{3mg}{k}$，选项 D 错误。

故选 BC。
}

\item
%% number: 9
%% paperName: 2025年甘肃高考第 9 题 5 分
%% typeId: 2
%% type: 多选题
%% chapter: 气体性质
%% point: 气体图象
%% method: 无
%% score: 5
%% degree: 550
%% duplicateId: 0
%% body:
如图，一定量的理想气体从状态 A 经等容过程到达状态 B，然后经等温过程到达状态 C。已知质量一定的某种理想气体的内能只与温度有关，且随温度升高而增大。下列说法正确的是\xzanswer{ACD}

\onepicture[w=4.30cm]{figs/09.png}

\fourchoices[answer=ACD]
{$A \to B$ 过程为吸热过程}
{$B \to C$ 过程为吸热过程}
{状态 A 压强比状态 B 的小}
{状态 A 内能比状态 C 的小}

%% answer: ACD
%% memo:
\memoanswer{
A．$A \to B$ 过程，体积不变，则 $W = 0$，温度升高，则 $\Delta U > 0$，根据热力学第一定律 $\Delta U = W + Q$ 可知 $Q > 0$，即该过程吸热，选项 A 正确；

B．$B \to C$ 过程，温度不变，则 $\Delta U = 0$，体积减小，则 $W > 0$，根据热力学第一定律 $\Delta U = W + Q$ 可知 $Q < 0$，即该过程为放热过程，选项 B 错误；

C．$A \to B$ 过程，体积不变，温度升高，根据查理定律可知，压强变大，即状态 A 压强比状态 B 压强小，选项 C 正确；

D．状态 A 的温度低于状态 C 的温度，可知状态 A 的内能比状态 C 的小，选项 D 正确。

故选 ACD。
}

\item
%% number: 10
%% paperName: 2025年甘肃高考第 10 题 5 分
%% typeId: 2
%% type: 多选题
%% chapter: 磁场
%% point: 带电粒子在磁场中的运动
%% method: 无
%% score: 5
%% degree: 350
%% duplicateId: 0
%% body:
2025 年 5 月 1 日，全球首个实现“聚变能发电演示”的紧凑型全超导托卡马克核聚变实验装置（BEST）在我国正式启动总装。如图是托卡马克环形容器中磁场截面的简化示意图，两个同心圆围成的环形区域内有垂直纸面向里的匀强磁场，磁感应强度大小为 $B$，内圆半径为 $R_{0}$。在内圆上 A 点有 a、b、c 三个粒子均在纸面内运动，并都恰好到达磁场外边界后返回。已知 a、b、c 带正电且比荷均为 $\frac{q}{m}$，a 粒子的速度大小为 $v_{a} = \frac{qBR_{0}}{m}$，方向沿同心圆的径向；b 和 c 粒子速度方向相反且与 a 粒子的速度方向垂直。不考虑带电粒子所受的重力和相互作用。下列说法正确的是\xzanswer{BD}

\onepicture[w=4.30cm]{\includesvg{figs/10}}

\fourchoices[answer=BD]
{外圆半径等于 $2R_{0}$}
{a 粒子返回 A 点所用的最短时间为 $\frac{(3\pi + 2)m}{qB}$}
{b、c 粒子返回 A 点所用的最短时间之比为 $\frac{\sqrt{2}}{\sqrt{2} + 2}$}
{c 粒子的速度大小为 $\frac{\sqrt{2}}{2}v_{a}$}

%% answer: BD
%% memo:
\memoanswer{
由题意，作出 a 粒子运动轨迹图，如图所示。

\onepicture[h=4.60cm]{figs/10_an1.png}

a 粒子恰好到达磁场外边界后返回，a 粒子运动的圆周正好与磁场外边界，然后沿径向做匀速直线运动，再做匀速圆周运动恰好回到 A 点，根据 a 粒子的速度大小为 $v_{a} = \frac{qBR_{0}}{m}$ 可得 $R_{a} = R_{0}$。

设外圆半径等于 $R'$，由几何关系得 $\angle AO'B = 270\Udu$

则 $R' = R_{0} + \sqrt{2}R_{0}$

A 错误；

B．由 A 项分析，a 粒子返回 A 点所用的最短时间为第一次回到 A 点的时间 $t_{\min}$

a 粒子做匀速圆周运动的周期 $T = \frac{2\pi R_{0}}{v_{a}} = \frac{2\pi m}{qB}$

在磁场中运动的时间 $t_{1} = \frac{540\Udu}{360\Udu} \cdot T = \frac{3\pi m}{qB}$

匀速直线运动的时间 $t_{2} = \frac{2R_{0}}{v_{a}} = \frac{2m}{qB}$

故 a 粒子返回 A 点所用的最短时间为 $t_{\min} = t_{1} + t_{2} = \frac{(3\pi + 2)m}{qB}$

B 正确；

C．由题意，作出 b、c 粒子运动轨迹图，如图所示

\twopicture[showlabel=false, wA=4.50cm, wB=4.50cm]{figs/10_an2.png}{figs/10_an3.png}

因为 b、c 粒子返回 A 点都是运动一个圆周，根据 b、c 带正电且比荷均为 $\frac{q}{m}$，所以两粒子做圆周运动周期相同，故所用的最短时间之比为 $1:1$，C 错误；

D．由几何关系得 $2R_{c} = \sqrt{2}R_{0}$

洛伦兹力提供向心力有 $qv_{c}B = \frac{mv_{c}^{2}}{R_{c}}$

联立解得 $v_{c} = \frac{\sqrt{2}}{2}v_{a}$

D 正确。

故选 BD。
}

\item
%% number: 11
%% paperName: 2025年甘肃高考第 11 题 6 分
%% typeId: 4
%% type: 实验题
%% chapter: 功和能
%% point: DIS研究机械能守恒定律
%% method: 无
%% score: 6
%% degree: 450
%% duplicateId: 0
%% body:
某学习小组使用如图~\ref{2025甘肃11a}所示的实验装置验证机械能守恒定律。

把一个直径为 $d$ 的小球用不可伸长的细线悬挂，光电门置于小球平衡位置处，其光线恰好通过小球球心，计时器与光电门相连。

将小球拉离平衡位置并记录其高度 $h$，然后由静止释放（运动平面与光电门光线垂直），记录小球经过光电门的挡光时间 $\Delta t$。改变 $h$，测量多组数据。已知重力加速度为 $g$，忽略阻力。

\begin{enumerate}
\item
以 $h$ 为横坐标、\tkanswer[2.20cm]{$\frac{1}{(\Delta t)^{2}}$}（填“$\Delta t$”、“$(\Delta t)^{2}$”、“$\frac{1}{\Delta t}$”或“$\frac{1}{(\Delta t)^{2}}$”）为纵坐标作直线图。若所得图像过原点，且斜率为\tkanswer[1.60cm]{$\frac{2g}{d^{2}}$}（用 $d$ 和 $g$ 表示），即可证明小球在运动过程中机械能守恒。
\item
实验中，用游标卡尺测得小球直径 $d = 20.48 \Umm$。
\begin{steps}
\item
由结果可知，所用的是\tkanswer[0.90cm]{50}分度的游标卡尺（填“10”、“20”或“50）；
\item
小组设计了一把 25 分度的游标卡尺，未测量时的状态如图~\ref{2025甘肃11b}所示。如果用此游标卡尺测量该小球直径、则游标尺上第\tkanswer[0.90cm]{12}条刻度线与主尺上的刻度线对齐。
\end{steps}
\end{enumerate}

\twopicture[labelA=2025甘肃11a, labelB=2025甘肃11b, align=right, wA=4.00cm, wB=9.00cm]{figs/11a.png}{figs/11b.png}

%% answer:
\jdanswer{
\begin{enumerate}
\item
$\frac{1}{(\Delta t)^{2}}$，$\frac{2g}{d^{2}}$
\item
①50　②12
\end{enumerate}
}
%% memo:
\memoanswer{
（1）小球经过光电门的挡光时间 $\Delta t$，可得小球到达平衡位置 $v = \frac{d}{\Delta t}$

为验证机械能守恒定律，此过程中重力势能转化为动能有 $mgh = \frac{1}{2}mv^{2}$

联立解得 $\frac{1}{(\Delta t)^{2}} = \frac{2g}{d^{2}}h$

可得纵坐标为 $\frac{1}{(\Delta t)^{2}}$

图像的斜率为 $k = \frac{2g}{d^{2}}$。

（2）10 分度、20 分度、50 分度的游标卡尺的精确度分别为 $0.1 \Umm$、$0.05 \Umm$、$0.02 \Umm$。此游标卡尺测得小球直径 $d = 20.48 \Umm$，可以判断所用的是 50 分度的游标卡尺。

若为 25 分度的游标卡尺，其精确度为 $0.04 \Umm$，用此游标卡尺测量该小球直径，可得

$n = \frac{0.48}{0.04} = 12$

则游标尺上第 12 条刻度线与主尺上的刻度线对齐。
}

\item
%% number: 12
%% paperName: 2025年甘肃高考第 12 题 9 分
%% typeId: 4
%% type: 实验题
%% chapter: 电路
%% point: 电路实验
%% method: 无
%% score: 9
%% degree: 450
%% duplicateId: 0
%% body:
某兴趣小组设计测量电阻阻值的实验方案。可用器材有：电池（电动势 $1.5 \UV$）两节，电压表（量程 $3 \UV$，内阻约 $3 \UkO$），电流表（量程 $0.3 \UA$，内阻约 $1 \UO$），滑动变阻器（最大阻值 $20 \UO$），待测电阻 $R_{x}$，开关 S$_{1}$，单刀双掷开关 S$_{2}$，导线若干。

\begin{enumerate}
\item
首先设计如图~\ref{2025甘肃12a}所示的电路。
\begin{steps}
\item
要求用 S$_{2}$ 选择电流表内、外接电路，请在图~\ref{2025甘肃12a}中补充连线将 S$_{2}$ 的 c、d 端接入电路\tkanswer[1.20cm]{如图}；
\item
闭合 S$_{1}$ 前，滑动变阻器的滑片 P 应置于\tkanswer[0.80cm]{b}端（填“a”或“b”）；
\item
闭合 S$_{1}$ 后，将 S$_{2}$ 分别接 c 和 d 端，观察到这两种情况下电压表的示数有变化、电流表的示数基本不变，因此测量电阻时 S$_{2}$ 应该接\tkanswer[0.80cm]{c}端（填“c”或“d”）。
\end{steps}
\item
为了消除上述实验中电表引入的误差、该小组又设计了如图~\ref{2025甘肃12b}所示的电路。
\begin{steps}
\item
请在图~\ref{2025甘肃12b}中补充连线将电压表接入电路\tkanswer[1.20cm]{如图}；
\item
闭合 S$_{1}$，将 S$_{2}$ 分别接 c 和 d 端时，电压表、电流表的读数分别为 $U_{c}$、$I_{c}$ 和 $U_{d}$、$I_{d}$。则待测电阻阻值 $R_{x} =$\tkanswer[2.60cm]{$\frac{U_{c}}{I_{c}} - \frac{U_{d}}{I_{d}}$}（用 $U_{c}$、$U_{d}$、$I_{c}$ 和 $I_{d}$ 表示）。
\end{steps}
\end{enumerate}

\twopicture[labelA=2025甘肃12a, labelB=2025甘肃12b, align=right, wA=7.60cm, wB=7.60cm]{figs/12a.png}{figs/12b.png}

%% answer:
\jdanswer{
\begin{enumerate}
\item
①如图　②b　③c
\item
①如图　②$\frac{U_{c}}{I_{c}} - \frac{U_{d}}{I_{d}}$
\end{enumerate}
}
%% memo:
\memoanswer{
（1）实物连接图如图所示。

\onepicture[w=7.50cm]{figs/12_an1.png}

闭合 S$_{1}$ 前，根据滑动变阻器的限流式接法，滑片 P 应置于 b 端，连入电路中的阻值最大，保护电路的安全。

闭合 S$_{1}$ 后，将 S$_{2}$ 分别接 c 和 d 端，观察到这两种情况下电压表的示数有变化、电流表的示数基本不变，说明电流表分压明显，为减小实验误差，应采用电流表外接法，因此测量电阻时 S$_{2}$ 应该接 c 端。

（2）实物连接图如图所示，

\onepicture[w=7.50cm]{figs/12_an2.png}

根据电路分析，当闭合 S$_{1}$，将 S$_{2}$ 接 c 端时，电压表、电流表的读数分别为 $U_{c}$、$I_{c}$，则

$\frac{U_{c}}{I_{c}} = R_{x} + R_{A} + R_{\text{滑}}$

将 S$_{2}$ 接 d 端时，电压表、电流表的读数分别为 $U_{d}$、$I_{d}$，则

$\frac{U_{d}}{I_{d}} = R_{A} + R_{\text{滑}}$

那么待测电阻阻值 $R_{x} = \frac{U_{c}}{I_{c}} - \frac{U_{d}}{I_{d}}$。
}

\item
%% number: 13
%% paperName: 2025年甘肃高考第 13 题 10 分
%% typeId: 6
%% type: 计算题
%% chapter: 光学
%% point: 光的折射
%% method: 无
%% score: 10
%% degree: 550
%% duplicateId: 0
%% body:
已知一圆台容器，高 $H = 15 \Ucm$，上口径 $R = 13 \Ucm$，容器底部中心有一质点，未装入水时，人眼从容器边缘无法观测到该质点，装入某种液体后，恰好可以看到，此时液面高度 $h = 12 \Ucm$，人眼观测角度 $\alpha$ 满足 $\sin \alpha = \frac{3}{5}$，人眼到入射处距离为 $5 \Ucm$。光在真空中的传播速度 $c = 3.0 \times 10^{8} \Ums$，求：

\begin{enumerate}
\item
该液体的折射率；
\item
光从底部质点反射至人眼全过程的时间。
\end{enumerate}

\onepicture[align=right, w=4.30cm]{figs/13.png}

%% answer:
\jdanswer{
\begin{enumerate}
\item
$\frac{4}{3}$
\item
$1 \times 10^{-9} \Us$
\end{enumerate}
}
%% memo:
\memoanswer{
（1）根据题意，画出光路图，如图所示。

\onepicture[w=4.30cm]{figs/13_an.png}

由几何关系可得 $\sin i = \frac{4}{5}$，$O'B = H - h = 3 \Ucm$

则有 $O'A = 4 \Ucm$，$AB = 5 \Ucm$

则 $\sin r = \frac{OC}{OB} = \frac{R - O'A}{\sqrt{h^{2} + (R - O'A)^{2}}} = \frac{3}{5}$

由折射定律可得该液体的折射率为 $n = \frac{\sin i}{\sin r} = \frac{4}{3}$

（2）根据题意，由图可知，光在空气中传播的距离为 $s_{1} = 10 \Ucm$

光在液体中的传播距离为 $s_{2} = OB = \sqrt{h^{2} + (R - O'A)^{2}} = 15 \Ucm$

光在液体中的传播速度为 $v = \frac{c}{n} = \frac{3c}{4}$

则光从底部质点反射至人眼全过程的时间 $t = \frac{s_{1}}{c} + \frac{s_{2}}{v} = \frac{0.1}{3 \times 10^{8}} \Us + \frac{0.15}{2.25 \times 10^{8}} \Us = 1 \times 10^{-9} \Us$
}

\item
%% number: 14
%% paperName: 2025年甘肃高考第 14 题 15 分
%% typeId: 6
%% type: 计算题
%% chapter: 运动和力
%% point: 牛顿定律的应用
%% method: 无
%% score: 15
%% degree: 450
%% duplicateId: 0
%% body:
如图~\ref{2025甘肃14a}所示，细杆两端固定，质量为 $m$ 的物块穿在细杆上。初始时刻。物块刚好能静止在细杆上。现以水平向左的力 $F$ 作用在物块上，$F$ 随时间 $t$ 的变化如图~\ref{2025甘肃14b}所示。开始滑动瞬间的滑动摩擦力等于最大静摩擦力。细杆足够长，重力加速度为 $g$，$\theta = 30\Udu$。求：

\begin{enumerate}
\item
$t = 6 \Us$ 时 $F$ 的大小，以及 $t$ 在 $0 \sim 6 \Us$ 内 $F$ 的冲量大小。
\item
$t$ 在 $0 \sim 6 \Us$ 内，摩擦力 $f$ 随时间 $t$ 变化的关系式，并作出相应的 $f-t$ 图像。
\item
$t = 6 \Us$ 时，物块的速度大小。
\end{enumerate}

\twopicture[labelA=2025甘肃14a, labelB=2025甘肃14b, numA=1, numB=2, align=right, wA=6.00cm, wB=5.40cm]{figs/14a.png}{figs/14b.png}

%% answer:
\jdanswer{
\begin{enumerate}
\item
$F = \frac{3\sqrt{3}mg}{2}$，$I = \frac{9\sqrt{3}mg}{2}$
\item
$f = \left( \frac{1}{2} - \frac{1}{8}t \right)mg$（$0 \leq t \leq 4$）；$f = \left( \frac{1}{8}t - \frac{1}{2} \right)mg$（$4 \leq t \leq 6$）
\item
$v = \frac{17}{2}g$
\end{enumerate}
}
%% memo:
\memoanswer{
（1）由图~\ref{2025甘肃14b}可知 $F$ 随时间线性变化，根据数学知识可知 $F = \frac{\sqrt{3}mg}{4}t$

所以当 $t = 6 \Us$ 时，$F = \frac{3\sqrt{3}mg}{2}$。

$0 \sim 6 \Us$ 内 $F$ 的冲量为 $F-t$ 图围成的面积，即 $I = \frac{1}{2} \times \frac{3\sqrt{3}}{2}mg \times 6 = \frac{9\sqrt{3}}{2}mg$

（2）由于初始时刻。物块刚好能静止在细杆上，则有 $mg\sin 30\Udu = \mu mg\cos 30\Udu$

即 $\mu = \tan 30\Udu = \frac{\sqrt{3}}{3}$

在垂直杆方向，当 $F\sin \theta = mg\cos \theta$ 时，$t = 4 \Us$

则 $0 \sim 4 \Us$，垂直杆方向 $F\sin \theta + N = mg\cos \theta$

摩擦力 $f = \mu N = \frac{\sqrt{3}}{3}\left( \frac{3}{2}mg - \frac{3}{8}mgt \right) = \left( \frac{1}{2} - \frac{1}{8}t \right)mg$（$0 \leq t \leq 4$）

在 $4 \sim 6 \Us$ 内，垂直杆方向 $F\sin \theta = mg\cos \theta + N$

摩擦力 $f = \mu N = \frac{\sqrt{3}}{3}\left( \frac{3}{8}mgt - \frac{3}{2}mg \right) = \left( \frac{1}{8}t - \frac{1}{2} \right)mg$（$4 \leq t \leq 6$）

相应的 $f-t$ 图像如图

\onepicture[w=5.00cm]{figs/14_an.png}

（3）在 $0 \sim 6 \Us$ 内沿杆方向根据动量定理有 $I_{F}\cos \theta - I_{f} + mg\sin \theta \times t = mv$

在 $0 \sim 6 \Us$ 内摩擦力的冲量为 $f-t$ 图围成的面积，则 $I_{f} = \frac{1}{2} \times \frac{1}{2}mg \times 4 + \frac{1}{2} \times \frac{1}{4}mg \times 2 = \frac{5}{4}mg$

联立有 $\frac{9\sqrt{3}}{2}mg \cdot \cos 30\Udu + 3mg - \frac{5}{4}mg = mv$

可得 $v = \frac{17}{2}g$
}

\item
%% number: 15
%% paperName: 2025年甘肃高考第 15 题 17 分
%% typeId: 6
%% type: 计算题
%% chapter: 电磁感应
%% point: 双杆问题
%% method: 微元
%% score: 17
%% degree: 350
%% duplicateId: 0
%% body:
在自动化装配车间，常采用电磁驱动的机械臂系统，如图，ab、cd 为两条足够长的光滑平行金属导轨，间距为 $L$，电阻忽略不计。导轨置于磁感应强度大小为 $B$，方向垂直纸面向里的匀强磁场中，导轨上有与之垂直并接触良好的金属机械臂 1 和 2，质量均为 $m$，电阻均为 $R$。导轨左侧接有电容为 $C$ 的电容器。初始时刻，机械臂 1 以初速度 $v_{0}$ 向右运动，机械臂 2 静止，运动过程中两机械臂不发生碰撞。系统达到稳定状态后，电流为零，两机械臂速度相同。

\begin{enumerate}
\item
求初始时刻机械臂 1 的感应电动势大小和感应电流方向；
\item
系统达到稳定状态前，若机械臂 1 和 2 中的电流分别为 $I_{1}$ 和 $I_{2}$，写出两机械臂各自所受安培力的大小；若电容器两端电压为 $U$，写出电容器电荷量的表达式；
\item
稳系统达到稳定状态后两机械臂的速度。若要两机械臂不相撞，二者在初始时刻的间距至少为多少？
\end{enumerate}

\onepicture[align=right, w=9.00cm]{figs/15.png}

%% answer:
\jdanswer{
\begin{enumerate}
\item
$BLv_{0}$，沿机械臂 1 向上
\item
$BI_{1}L$，$BI_{2}L$，$Q = CU$
\item
$\frac{mv_{0}}{CB^{2}L^{2} + 2m}$，方向向右；$\frac{mv_{0}R}{B^{2}L^{2}}$
\end{enumerate}
}
%% memo:
\memoanswer{
（1）由法拉第电磁感应定律可知，初始时刻机械臂 1 的感应电动势大小为 $E = BLv_{0}$

由右手定则可知感应电流方向沿机械臂 1 向上。

（2）在达到稳定前，两机械臂电流分别为 $I_{1}$ 和 $I_{2}$，两机械臂安培力的大小分别为 $F_{1} = BI_{1}L$，$F_{2} = BI_{2}L$。

设电容器所带电荷量为 $Q$，则 $Q = CU$。

（3）达到稳定 $I_{1} = I = 0$ 时，两机械臂的速度相同，产生的感应电动势与电容器的电压相等，回路中没有电流。结合（2）问的分析可知此时 $I_{1} = I_{2} = 0$，$U = \frac{BLmv_{0}}{2m + B^{2}L^{2}C}$

同时 $U = BLv$

可得两机械臂的速度为 $v = \frac{mv_{0}}{CB^{2}L^{2} + 2m}$

方向向右

结合（2）问分析，在任意时刻有 $U = BLv_{1} - I_{1}R = BLv_{2} + I_{2}R$

即 $BL(v_{1} - v_{2}) = I_{2}R + I_{1}R$

对该式两边取全过程时间的累计有 $BL(v_{1} - v_{2})\Delta t = I_{2}R\Delta t + I_{1}R\Delta t$

其中 $(v_{1} - v_{2})\Delta t = x_{1} - x_{2} = d_{\min}$，$I_{1}\Delta t = Q_{1}'$，$I_{2}\Delta t = Q_{2}'$

即 $BLd_{\min} = (Q_{1}' + Q_{2}')R$

从开始到最终稳定的过程中，对机械臂 1 和机械臂 2 分别根据动量定理有

$-B\overline{I}_{1}'L \cdot \Delta t = mv - mv_{0}$，$B\overline{I}_{2}'L \cdot \Delta t = mv$

即 $BLQ_{1}' = mv_{0} - mv$，$BLQ_{2}' = mv$

可得 $BL(Q_{1}' + Q_{2}') = mv_{0}$

联立解得稳定时的速度和两棒间初始距离的最小值为 $d_{\min} = \frac{mv_{0}R}{B^{2}L^{2}}$。
}

\end{enumerate}

\end{document}
